\documentclass[12pt,a4paper]{book}
\usepackage[utf8]{inputenc}
\usepackage[T1]{fontenc}
\usepackage[french]{babel}
\usepackage{amsmath, amsthm, amssymb}
\usepackage{graphicx}
\usepackage{hyperref}
\usepackage{geometry}
\geometry{a4paper, margin=1in}

% Définitions des environnements mathématiques
\newtheorem{theorem}{Théorème}[chapter]
\newtheorem{definition}{Définition}[chapter]
\newtheorem{lemma}{Lemme}[chapter]
\newtheorem{corollary}{Corollaire}[chapter]

\title{\textbf{Maths for AI}\\Les Fondements Mathématiques de l'Intelligence Artificielle}
\author{Instructeur}
\date{\today}

\begin{document}

\maketitle
\tableofcontents

\chapter{Algèbre Linéaire \& Tenseurs}

L'algèbre linéaire est le langage universel dans lequel sont écrites les données en intelligence artificielle. Des simples vecteurs aux tenseurs multidimensionnels, comprendre ces structures est essentiel pour manipuler l'information.

\section{Espaces Vectoriels}

\begin{definition}[Espace Vectoriel]
Un espace vectoriel $V$ sur un corps $K$ (généralement $\mathbb{R}$ ou $\mathbb{C}$) est un ensemble muni de deux lois (addition interne et multiplication par un scalaire) vérifiant des propriétés de commutativité, d'associativité et de distributivité.
\end{definition}

\section{Matrices et Applications Linéaires}

Les matrices permettent de représenter les applications linéaires entre espaces vectoriels de dimension finie. Dans les réseaux de neurones, les poids d'une couche dense (Fully Connected) sont représentés par une matrice $W$.

\begin{theorem}[Décomposition en Valeurs Singulières - SVD]
Toute matrice réelle $A \in \mathbb{R}^{m \times n}$ peut être décomposée sous la forme :
$$ A = U \Sigma V^T $$
où $U$ est une matrice orthogonale $m \times m$, $V$ est une matrice orthogonale $n \times n$, et $\Sigma$ est une matrice rectangulaire diagonale $m \times n$ contenant les valeurs singulières de $A$.
\end{theorem}

\begin{proof}
La preuve repose sur le fait que $A^TA$ et $AA^T$ sont des matrices symétriques réelles et semi-définies positives. Leurs vecteurs propres forment les colonnes de $V$ et $U$ respectivement, tandis que les valeurs singulières sont les racines carrées de leurs valeurs propres (qui sont communes et réelles positives).
\end{proof}

\section{Calcul Tensoriel pour le Deep Learning}

En Deep Learning, les tenseurs généralisent le concept de vecteurs et de matrices à des dimensions supérieures. Un vecteur est un tenseur d'ordre 1, une matrice d'ordre 2, et une image RGB d'ordre 3.

\chapter{Optimisation}

L'optimisation est au cœur de l'apprentissage automatique : on cherche en permanence à minimiser une fonction de coût (ou perte) par rapport aux paramètres d'un modèle.

\section{Convexité}

\begin{definition}[Fonction Convexe]
Une fonction $f : \mathbb{R}^n \rightarrow \mathbb{R}$ est convexe si, pour tous $x, y \in \mathbb{R}^n$ et pour tout $\theta \in [0, 1]$, on a :
$$ f(\theta x + (1-\theta)y) \leq \theta f(x) + (1-\theta)f(y) $$
\end{definition}

\begin{theorem}[Minimum Global]
Si une fonction de coût $J$ est strictement convexe, alors tout minimum local de $J$ est un minimum global unique.
\end{theorem}

\section{Descente de Gradient}

La descente de gradient est l'algorithme d'optimisation du premier ordre le plus utilisé pour minimiser une fonction de coût différentiable $J(\theta)$.

\begin{definition}[Mise à jour du gradient]
L'itération de la descente de gradient est donnée par :
$$ \theta_{t+1} = \theta_t - \alpha \nabla J(\theta_t) $$
où $\alpha$ est le taux d'apprentissage (learning rate) et $\nabla J(\theta_t)$ est le gradient de la fonction de coût au point $\theta_t$.
\end{definition}

\section{Multiplicateurs de Lagrange}

Pour les problèmes d'optimisation sous contraintes, on utilise la méthode des multiplicateurs de Lagrange.

\begin{definition}[Lagrangien]
Pour minimiser $f(x)$ sous la contrainte d'égalité $g(x) = 0$, on introduit le Lagrangien :
$$ \mathcal{L}(x, \lambda) = f(x) + \lambda g(x) $$
Les points critiques sont trouvés en résolvant $\nabla \mathcal{L}(x, \lambda) = 0$.
\end{definition}

\chapter{Probabilités \& Information}

L'IA ne modélise pas la certitude absolue, mais navigue dans l'incertitude. La théorie des probabilités offre le cadre mathématique pour gérer cette incertitude, tandis que la théorie de l'information quantifie les données.

\section{Distributions de Probabilités}

\begin{definition}[Distribution Normale]
La loi normale (ou Gaussienne) de moyenne $\mu$ et d'écart-type $\sigma$ a pour densité de probabilité :
$$ f(x) = \frac{1}{\sigma\sqrt{2\pi}} e^{-\frac{1}{2}\left(\frac{x-\mu}{\sigma}\right)^2} $$
\end{definition}

\section{Théorie de l'Information}

\begin{definition}[Entropie de Shannon]
L'entropie d'une variable aléatoire discrète $X$ ayant une distribution de probabilité $P$ est définie par :
$$ H(X) = - \sum_{x \in X} P(x) \log P(x) $$
Elle mesure la quantité d'incertitude ou de "surprise" inhérente aux résultats possibles de $X$.
\end{definition}

\section{Entropie Croisée}

L'entropie croisée est la fonction de perte la plus courante pour les problèmes de classification en apprentissage automatique.

\begin{definition}[Entropie Croisée]
L'entropie croisée entre une vraie distribution $P$ et une distribution estimée $Q$ est :
$$ H(P, Q) = - \sum_{x \in X} P(x) \log Q(x) $$
Minimiser l'entropie croisée équivaut à minimiser la divergence de Kullback-Leibler (KL) entre $P$ et $Q$.
\end{definition}

\chapter{Calcul Différentiel}

L'apprentissage des réseaux de neurones repose sur l'ajustement de leurs paramètres en fonction de l'erreur produite. Cette adaptation nécessite de calculer des dérivées pour propager l'erreur.

\section{Dérivées Partielles et Gradient}

\begin{definition}[Gradient]
Soit une fonction $f : \mathbb{R}^n \rightarrow \mathbb{R}$. Le gradient de $f$, noté $\nabla f(x)$, est le vecteur des dérivées partielles par rapport à chaque coordonnée :
$$ \nabla f(x) = \left( \frac{\partial f}{\partial x_1}, \frac{\partial f}{\partial x_2}, \dots, \frac{\partial f}{\partial x_n} \right)^T $$
\end{definition}

\section{Jacobiennes et Hessiennes}

\begin{definition}[Matrice Jacobienne]
Pour une fonction vectorielle $f : \mathbb{R}^n \rightarrow \mathbb{R}^m$, la Jacobienne $J$ est la matrice $m \times n$ dont l'élément $(i,j)$ est $\frac{\partial f_i}{\partial x_j}$.
\end{definition}

\begin{definition}[Matrice Hessienne]
Pour une fonction scalaire $f : \mathbb{R}^n \rightarrow \mathbb{R}$, la matrice Hessienne $H$ est la matrice carrée $n \times n$ des dérivées partielles secondes :
$$ H_{i,j} = \frac{\partial^2 f}{\partial x_i \partial x_j} $$
\end{definition}

\section{Rétropropagation (Backpropagation)}

\begin{theorem}[Règle de la Chaîne (Chain Rule)]
Si $y = g(x)$ et $z = f(y)$, alors la dérivée de $z$ par rapport à $x$ est donnée par le produit des Jacobiennes :
$$ \frac{\partial z}{\partial x} = \frac{\partial z}{\partial y} \frac{\partial y}{\partial x} $$
\end{theorem}

\begin{proof}
La rétropropagation exploite directement cette règle de la chaîne de manière récursive, en partant de la fonction de perte (la sortie) et en remontant couche par couche jusqu'à l'entrée, ce qui permet un calcul efficace du gradient par rapport à tous les poids du réseau.
\end{proof}

\chapter{Géométrie des Données}

Les données modernes (images, textes, graphes) résident souvent dans des espaces de très haute dimension. Cependant, elles sont souvent contraintes à vivre sur des sous-espaces de plus faible dimension appelés variétés (manifolds).

\section{Variétés (Manifolds)}

L'hypothèse des variétés stipule que les données réelles de grande dimension (par exemple, des images de visages) se concentrent près d'une variété de faible dimension plongée dans l'espace de grande dimension.

\begin{definition}[Variété Différentielle]
Une variété différentielle de dimension $d$ est un espace topologique qui ressemble localement à l'espace euclidien $\mathbb{R}^d$.
\end{definition}

\section{Espaces Latents}

En apprentissage profond (notamment avec les Auto-encodeurs et les GANs), le réseau apprend à projeter les données depuis l'espace d'entrée (de haute dimension) vers un espace latent de dimension réduite, en capturant les caractéristiques sémantiques.

\section{Métriques de Distance}

Dans ces espaces latents, la distance euclidienne standard n'est souvent pas appropriée pour mesurer la similarité sémantique.

\begin{definition}[Distance de Minkowski]
La distance de Minkowski d'ordre $p$ entre deux vecteurs $x, y \in \mathbb{R}^n$ est définie par :
$$ D_p(x, y) = \left( \sum_{i=1}^n |x_i - y_i|^p \right)^{1/p} $$
\end{definition}

\begin{corollary}[Cas particuliers]
Pour $p=1$, c'est la distance de Manhattan. Pour $p=2$, c'est la distance euclidienne classique.
\end{corollary}


\end{document}
